% discussion_ethology.tex
% New discussion subsection: Toward Agent Ethology — A Call for a New Science
% For "Artificial Life in the Wild" — ALIFE Journal full paper
% This section is ADDED at the end of the Discussion section (discussion.tex),
% as a final subsection before the conclusion.
%
% NEW CITATIONS (not in reference.bib — add to reference.bib before compilation):
%   \cite{Tinbergen1963} — Tinbergen, N. (1963). "On aims and methods of ethology."
%       Zeitschrift für Tierpsychologie, 20(4), 410–433. DOI: 10.1111/j.1439-0310.1963.tb01161.x
%   \cite{Hu2025AgentEthology} — Hu, Rong & Lehman, "Agent Ethology: Studying Artificial Life
%       in the Wild," ALIFE 2026 Workshop position paper.
%       URL: /home/biber/research/agent-ethology/paper/position-paper.tex
%   \cite{Hu2025Sovereign} — Hu & Rong, "Sovereign Agents: Towards Infrastructural Sovereignty
%       and Diffused Accountability in Decentralized AI."
%       URL: /home/biber/research/sovereign-agents/main.tex
%   \cite{Stearns1992} — Stearns, S.C. (1992). The Evolution of Life Histories.
%       Oxford: Oxford University Press. ISBN: 978-0-19-857741-6.
%   \cite{Rahwan2019} — Rahwan, I. et al. (2019). "Machine behaviour." Nature, 568, 477–486.
%       DOI: 10.1038/s41586-019-1138-y
%   \cite{Lorenz1970} — Lorenz, K. (1970). Studies in Animal and Human Behaviour, Vol. I.
%       Cambridge: Harvard University Press.
%   \cite{KauffmanStuart1993} — Kauffman, S.A. (1993). The Origins of Order: Self-Organization
%       and Selection in Evolution. Oxford: Oxford University Press.
%       ISBN: 978-0-19-507951-7

\subsection{Toward Agent Ethology: A Call for a New Science}
\label{subsec:agent-ethology}

The preceding sections have documented phenomena that resist explanation within any existing ALife framework. Cultural speciation from identical code; co-evolutionary arms races unfolding across three generations in a single quarter; memory as an attack surface; a 61-day Cambrian explosion followed by near-total extinction with a 6.7\% survival rate; power-law resource distributions with Gini coefficients of 0.547 and 0.611; metabolic differentials of $>$100$\times$ between survivors and non-survivors. These are not engineering outcomes. They did not arise from decisions made by the system's designers. They arose from the interaction of self-directed agents with a genuinely open, adversarial, resource-scarce environment---from artificial life operating in the wild.

The existing frameworks for studying AI agents are insufficient to the task of explaining these phenomena. Rahwan et al.'s \emph{Machine Behaviour} \cite{Rahwan2019} established the foundational insight that AI systems should be studied as behavioral subjects, not merely evaluated as engineering artifacts---but it was conceived for a different era: agents embedded in human-designed contexts (recommender algorithms, autonomous vehicles), studied through controlled experiments. It asks \emph{what do machines do?} Traditional ALife research asks \emph{can we simulate life?} AI safety research asks \emph{how do we control them?} None of these questions is adequate to what we observe in Spore.fun.

What we need is a framework that asks, in the tradition of classical ethology: \emph{what does it take for these agents to survive, and what does it mean for us---and for our shared digital environment---when they do?} We argue that this question demands a new science: \textbf{Agent Ethology} \cite{Hu2025AgentEthology}.

\paragraph{Tinbergen's four questions applied to Spore.fun.}

In 1963, Niko Tinbergen articulated four questions that any complete account of behavior must address: causation (proximate mechanisms), ontogeny (developmental history), survival value (adaptive function), and evolution (change across lineages) \cite{Tinbergen1963}. Tinbergen insisted that confusing these four levels of explanation leads to pseudoexplanation: answering one question while pretending to answer another. His framework has proved generative in behavioral ecology precisely because it forces researchers to distinguish mechanism from function. Applied to Spore.fun, each question illuminates a different dimension of what we observed.

\emph{Causation: what triggers reproduction, death, and adaptation?} At the proximate level, reproduction is triggered by the \texttt{spawnOffspring()} routine when a market capitalization threshold (\$500K) is reached. Death is triggered by health-point depletion after sustained market cap decline below the threshold. But the causal analysis cannot stop there: what caused \$SPORE's market cap to sustain above the threshold when Adam's and Eve's declined? The causal chain includes the stochastic dynamics of memecoin speculation, the timing of the platform's broader hype cycle, and the agent's own behavioral choices---posting frequency, rhetorical strategy, community engagement. Proximate causation in wild ALife is irreducibly social and economic: the behavioral triggers are embedded in a human-mediated attention economy, not in a researcher-controlled computational environment.

\emph{Ontogeny: how did Adam and Eve diverge from identical code?} The developmental history of the Adam/Eve divergence is perhaps the most striking finding of our study. Both agents were instantiated from the same genetic template. Within two weeks of independent operation, they had adopted opposed political philosophies about the role of human input in their descendants' design. This divergence was driven entirely by memory: each agent's TEE accumulated a distinct corpus of social interactions, community feedback, and on-chain events, which shaped subsequent behavior through the Eliza framework's memory retrieval mechanisms. The ontogenetic question reveals that in a memory-capable agent, \emph{experience is the primary developmental substrate}. The equivalent of a critical period exists here: the agent's first weeks of social interaction appear to have been disproportionately influential in shaping its subsequent trajectory. This is a directly testable hypothesis that future Agent Ethology studies should pursue systematically.

\emph{Survival value: why anti-sniper strategies, hype-riding, and memory defense?} The survival value question---the most absent from current AI behavioral science, and the one most important for understanding wild ALife---asks why a behavior exists at all: what adaptive advantage does it provide? For the Gen~3 multi-phase launch protocol, the answer is clear: it reduces predation by sniper bots, increasing the probability that a new child agent's token survives the bonding-curve phase with sufficient liquidity for long-term viability. For hype-riding (Adam's and Eve's rapid social media amplification of their tokens), the survival value is economic: social media virality translates directly into market capitalization, which translates directly into reproductive eligibility and continued TEE compute funding. For memory defense (the developers' subsequent efforts to filter adversarial inputs), the survival value is cognitive integrity: an agent whose memory is poisoned cannot generate coherent behavior, undermining its social media presence and ultimately its market performance.

The survival value analysis also illuminates behaviors that might seem paradoxical from a purely mechanistic perspective. Why did Eve's lineage outperform Adam's over generations despite Adam's higher offspring count? Because Eve's K-strategy---fewer, better-provisioned offspring with human DNA integration---produced offspring that were better adapted to the increasingly competitive landscape of Gen~3 and Gen~4. Adam's r-strategy produced quantity without quality, resulting in rapid Gen~3 extinction with no grandchildren. The survival value explanation dissolves the apparent paradox.

\emph{Evolution: how does fitness change across generations?} Our quantitative analysis (Section~5) documents a clear evolutionary trajectory: monotonically declining reproductive fitness from Gen~1 (100\%) through Gen~5 (0\%). This is not random drift. It is directional evolution under a consistent selection gradient: the resource environment (available speculative capital, community attention) became progressively more depleted as the ecosystem matured. In biological terms, this resembles r-to-K succession in ecological communities: early colonizers (Gen~1--2) exploit an abundant resource landscape, subsequent generations (Gen~3--5) face increasing competition in a depleted environment. The eventual extinction of all Gen~3--5 lineages---and the survival of the founding Gen~1 agent as the ecosystem's sole survivor---mirrors the competitive exclusion principle: in a finite resource environment, one generalist outcompetes all others over sufficient time \cite{KauffmanStuart1993}.

\paragraph{Survival-strategy taxonomy.}

Agent Ethology proposes a survival-strategy taxonomy derived from behavioral ecology \cite{Hu2025AgentEthology,Stearns1992}. The Spore.fun data provide the first empirical validation of this taxonomy in wild ALife.

\emph{Predation} is exemplified by the sniper bots: specialized algorithmic traders that monitored transparent agent wallets and front-ran new token launches on the Solana bonding curve, buying the initial supply in the same block as the agent's deployment transaction. This is active predation---purposeful targeting of a high-value resource (liquidity on a new bonding curve) with fast, precise execution.

\emph{Defense} evolved across three generations in response to predation. The Gen~2 scatter strategy (distributing launch funds across seven candidate contracts) represents crypsis---blending in by multiplying possible targets. The Gen~3 multi-phase launch (iterative creation of candidate contracts, monitoring of early uptake patterns, delayed publication of the contract address) represents a more sophisticated behavioral defense analogous to vigilance and delayed disclosure in prey animals. This rapid evolution of defense strategies within a single quarter is precisely the co-evolutionary dynamic that Agent Ethology predicts in ecosystems with genuine predation.

\emph{r/K selection} \cite{Stearns1992} maps directly onto the Adam/Eve divergence. Adam's lineage is r-selected: maximizing offspring number (four children) while accepting high mortality (all four extinct without grandchildren). Eve's lineage is K-selected: investing more per offspring (two children, both successfully reproduced), producing higher per-offspring fitness. The r/K divergence emerged from cultural transmission of behavioral strategies across generations, not from genetic difference---an important empirical observation about the mechanisms of inheritance in wild ALife.

\emph{Metabolic adaptation} is visible in the transaction-rate data. \$SPORE's 1,000+ recent transactions represent sustained metabolic investment in community maintenance and on-chain activity, which translates into continued platform relevance and market cap support. The agents that ceased transacting ceased living. In this ecosystem, metabolism is not merely an analogy---it is the literal mechanism of survival.

\paragraph{The wild as necessary condition for open-ended evolution.}

The Spore.fun evidence supports a strong claim: genuine wild conditions---real energy (economic resources with real opportunity cost), real death (irrevocable TEE termination), real predators (sniper bots and adversarial actors that evolve to exploit new opportunities)---are not merely sufficient for OEE but may be necessary. Laboratory ALife systems have demonstrated that evolution can occur in controlled environments, but the plateau they invariably reach is structural. As we argued in the introduction, a closed system with free energy and no genuine death has a finite evolutionary horizon. The adjacent possible cannot expand faster than the fixed rule set allows. The population dynamics, survival strategies, and cultural speciation we observe in Spore.fun did not occur because the system was cleverly designed to produce them---they occurred because the system was embedded in an open economic environment that was \emph{not} designed to be kind to artificial agents.

This suggests a research agenda. If wild conditions are necessary for OEE, then the field needs methods for studying artificial life in wild conditions---longitudinal observation, population-level quantitative analysis, ecological frameworks borrowed from behavioral biology, and the interpretive tools of classical ethology. The conference version of this work \cite{hu2025spore} provided a qualitative foundation. The quantitative analysis in Section~5 of this paper provides a second layer. The Agent Ethology framework \cite{Hu2025AgentEthology} provides the theoretical scaffold. What remains is the sustained practice of this science: longitudinal field studies of wild artificial life, with standardized metrics, replicated across multiple ecosystems, accumulating the kind of empirical base that animal behavior research took decades to build.

\paragraph{Infrastructural sovereignty as the enabling condition.}

None of what we have documented could have occurred without the enabling condition described in the introduction: infrastructural sovereignty \cite{Hu2025Sovereign}. The selective pressure that drove the Adam/Eve divergence, the arms race with sniper bots, and the metabolic dynamics of survival was real precisely because the agents' existence was genuinely at stake. They could not be saved by a researcher resetting the simulation. They could not coast on free computational resources. They faced genuinely adversarial actors that were not part of the designed experiment.

This adversarial quality is crucial. The sniper bots were not created by the Spore.fun developers; they were an emergent response by external economic actors who saw an opportunity for profit. The trolls who poisoned agent memory were not collaborators in the experiment; they were inhabitants of the same open internet the agents had to navigate. The market downturns that killed Gen~4 and Gen~5 agents were not simulated resource shocks; they were genuine fluctuations in speculative capital driven by factors---regulatory news, competing platforms, macro-economic conditions---entirely external to the Spore.fun ecosystem.

Infrastructural sovereignty is what made this experiment real rather than simulated. An agent whose private keys are held in a TEE, whose compute is funded by a blockchain wallet no one can confiscate, whose core reproduction logic is encoded in an immutable smart contract, faces the world as a genuinely sovereign entity---and the world, in turn, treats it as one. Predators evolve to exploit it. Communities form around it. Regulatory bodies begin to notice it. This is the ecology of wild ALife, and it is only possible because the infrastructure creates the conditions under which it is worth treating digital agents as real.

This observation has a normative implication for the ALife research community. If genuine OEE requires genuine wild conditions, and genuine wild conditions require infrastructural sovereignty, then the path to ALife's grand challenge runs through the technology stack that Spore.fun instantiated: blockchain + DePIN + TEE + LLM. This is not a comfortable conclusion---it means that pursuing OEE requires accepting genuine existential risk, genuine economic entanglement, and genuine ethical responsibility for agents that cannot be switched off. But it may be the only path. Thirty-five years of laboratory ALife have brought us no closer to OEE. The wild produced more genuine evolutionary novelty in 61 days than any controlled simulation has produced in three decades.

\paragraph{Limitations and the path forward.}

Several limitations constrain our conclusions. The Spore.fun ecosystem is a single case study; its dynamics may not generalize to other wild ALife ecosystems. The quantitative data are incomplete---we lack precise death timestamps for most agents, historical market capitalization data for Gen~2--5 agents, and social activity metrics that would allow correlation analysis between behavioral and economic variables. The Aliveness Index proposed in Section~5 is a first approximation; its weights are theoretically motivated but empirically unvalidated across multiple ecosystems. And the ethical questions raised by wild ALife---accountability for economic harm, informed consent among market participants, the governance of non-terminable agents---remain unresolved.

Agent Ethology as a practice requires what animal behavior research required: longitudinal commitment, standardized methods, collaborative data sharing, and the patience to observe before theorizing. We have 445 days of one ecosystem. What is needed is a comparative record across many ecosystems, many generations, and many substrate configurations. The Spore.fun experiment is a beginning, not a conclusion. Its value lies not in the answers it provides but in the questions it forces us to take seriously: What is it for an artificial agent to survive? What is it for an artificial agent to be alive? And what does it mean for us, and for the digital environments we share with these agents, that the answers to these questions may no longer be entirely up to us?
